\section{Related Work}\label{sec:related}
%==============================================================================

\subsection{Zero-Error and Side-Information Lineage}\label{sec:related-source-coding}

Shannon's zero-error framework and its graph-theoretic refinements by Körner and Lovász provide the closest conceptual starting point~\cite{shannon1956zero,korner1973graphs,lovasz1979shannon}. Witsenhausen's zero-error side-information problem is directly relevant because it studies exact recovery under side information through graph coloring and label budgets~\cite{witsenhausen1976zero}. Our one-shot coloring statements lie in that lineage. As in Witsenhausen's formulation, an observation law induces the relevant confusability graph. The difference is that here the observation law is generated by a deterministic multi-fact tuple-space architecture with admissible coordinate-subset views, so the model itself constrains what graphs can arise.

\paragraph{Key distinction from Witsenhausen.}
Both settings can be viewed as starting from an observation law and passing to the induced confusability graph, so there is no sharp ``induced'' versus ``given'' divide. The difference is architectural and structural. Witsenhausen's setting starts from a side-information law and studies the coding problem for the graph it induces. Our setting starts from a deterministic partial-view architecture on latent tuples, with observations obtained by admissible coordinate-subset projections, and asks what graph structure that architecture can generate and what recovery laws follow from it. This makes two model-side issues central: how the admissible view family constrains the induced graph class, and how changing the view topology changes the resulting failure graph and recovery budget. In the exact full-tuple-space coordinate-view model analyzed here, that graph class can now be characterized exactly: the realizable confusability relations are precisely those determined by upward-closed families of coordinate-agreement sets. The contribution is therefore not a new coloring theorem for arbitrary graphs, but an architectural realization theory connecting multi-fact tuple-space views, a fully characterized model-generated graph class, and exact zero-error recovery.

Slepian--Wolf coding and classical entropy converses supply the second line of influence~\cite{slepian1973noiseless,fano1961transmission,witsenhausen1975conditional,cover2006elements}. Coding-for-computing, functional compression, and graph-entropy/characteristic-graph viewpoints are also close neighbors because they study deterministic decoder targets and coding questions once an underlying graph or function structure has been specified~\cite{orlitsky2001coding,doshi2010functional,korner1973graphs}.

\medskip
\noindent\textbf{Characteristic‑graph vs. coordinate‑projection models.} Characteristic‑graph and functional‑compression formalisms handle arbitrary deterministic functions of the source and therefore subsume a broader class of recovery tasks than the present model. Our model intentionally restricts attention to coordinate‑projection observations on labeled tuple spaces rather than arbitrary functions. That restriction is the source of our structural leverage: it allows a reduction of confusability to agreement‑set membership and thereby enables proofs that (a) realizable confusability relations are exactly those determined by upward‑closed agreement families, (b) coordinatewise alphabet permutations act by graph automorphisms, and (c) the induced graph class is monotone and closed under the block‑composition (strong‑product) operation. These specific structural consequences do not follow from the general characteristic‑graph viewpoint and are the reasons the paper can give deductive, architecture‑first results rather than only applying general functional‑compression theorems.

In our setting the side information is deterministic rather than stochastic, but it plays the same operational role: exact recovery becomes impossible when the available observation/tag alphabet is too small to separate the remaining alternatives.

After a confusability graph is fixed, the colorability, strong-power, Shannon-capacity, and Lov\'asz-$\vartheta$ machinery used later is classical. In particular, capacity--$\vartheta$ equality for cluster graphs is classical; the novelty here is identifying transitivity of view-generated confusability as the model-side condition that produces this equality subclass, together with meet-witnessing and fiber coherence as structural sufficient conditions. The paper therefore does not claim a new theorem about arbitrary graphs, but a deterministic partial-view model that generates a fully characterized monotone agreement-set graph class on the full labeled tuple space, an exact recovery/success-set/weighted-success theory for that model-generated class, and a structural equality route inside it. What it does not attempt is the broader unlabeled-isomorphism classification, or the corresponding classifications for restricted realized-state families and stochastic support-overlap models.

\subsection{Consistency-Constrained Storage and Interaction}\label{sec:related-distributed}

Consistency-constrained storage problems, including multi-version coding and related distributed-storage converses, show that correctness requirements impose unavoidable information costs~\cite{rashmi2015multiversion}. Here the object under study is a modifiable encoding architecture, and the main question is the exact-consistency cost of multiple independently writable representations of one latent fact.

\subsection{Computational Realizations}\label{sec:related-meta}

Reflection, metadata systems, and maintenance mechanisms in host platforms provide examples of how the boundary-case realizability conditions can be instantiated in practice~\cite{kiczales1991art,smith1984reflection}. These examples are secondary to the paper's main zero-error and graph-capacity contribution.

\subsection{Formal Verification}\label{sec:related-formal}

The Lean 4 formalization places the work in the tradition of mechanized mathematics and verified theory development~\cite{demoura2021lean4,leroy2009compcert}.

%==============================================================================
